\chapter{Auction Mechanics and Theory}\label{ch:mx:auction-mechanics-and-theory}

In the minutes before a closing auction the exchange publishes an indicative price and imbalance every few seconds, and much of the auction's volume
arrives in its last seconds. Every participant watches the same numbers and decides when to show its hand. The closing auction sets the prices that
index funds track and that funds are valued at, so its design, what is published, when orders are accepted and when exactly it ends, decides who can
move the close and who pays for it. This chapter runs closing calls in the exchange simulator, measures how the indicative converges to the final
price, prices a late order under a fixed and a \omterm{def:mx:auction-mechanics-and-theory:randomend}{random end}, and replays the sniping race that \omterm{def:mx:auction-mechanics-and-theory:fba}{frequent batch auctions} were proposed to stop.

\section{Uncrossing rules revisited}

A call auction (One Quant Book 1, chapter~13) collects orders without matching them and clears them all at one price: the price that maximises the
executed volume, then minimises the surplus, then follows the side with the surplus, then stays closest to a reference price (Book~1's
\texttt{firm.auction}). With buys of 300 shares at market, 200 at 10.02 and 400 at 10.00, and sells of 500 at 10.00, 300 at 10.01 and 200 at 10.03, the volume
is 500 at 10.00, 10.01 and 10.02; the surplus is 400 to buy at 10.00 and 300 to sell at the other two, and the sell pressure picks the lower, 10.01.
Madhavan (1992) compared the mechanisms: a periodic auction offers greater price efficiency and can function where continuous mechanisms fail, but traders
give up continuity and bear higher information costs.

\texttt{firm.exchsim}'s engine runs the calls: orders are accepted without matching, market-on-close and limit-on-close orders (tif C) are kept aside, the
engine publishes an indicative price and imbalance on a schedule (feed message I), and the uncross calls \texttt{firm.auction}. \texttt{firm.auctionsim}
states a design as a policy and turns it into the simulator's configuration.

\omcode{code/firm/auctionsim/firm_auctionsim.py}{45}{57}{A closing auction as a policy: the call's phases, and a collar as a price band from the start of the call.}\label{lst:mx:auction-mechanics-and-theory:policy}

\section{Indicative prices and imbalance publication}

The indicative is the price the auction would clear at if it ended now, and the imbalance the shares that would be left on one side. Publishing them
invites the other side: a large buy imbalance at a price above value is an offer to whoever will sell. The chapter's closing call lasts five minutes with
an indicative every second. Standing limit-on-close interest of 500 shares a tick sits on each side of 100.00; random market-on-close orders of 100 to 1\,000
shares arrive five times a minute until five seconds before the end; a fund sends a market-on-close buy of 10\,000 shares a minute into the call.
Responders read each indicative and, half the time, when the price has been pushed more than a tick from 100.00, offer the other side at 100.01 (or 99.99)
in proportion to the distance, 250 shares a tick up to 3\,000: they supply the imbalance at the price they believe, not at the price it has been pushed to.

\omcode{code/microstructure/10-auction-mechanics-and-theory/python/mx_auction.py}{57}{68}{A responder to the published imbalance: it offers the other side at its own value plus a tick, in proportion to the push.}\label{lst:mx:auction-mechanics-and-theory:resp}

\begin{omfigure}
\begin{tikzpicture}
\begin{groupplot}[group style={group size=2 by 1,horizontal sep=1.6cm},width=6.6cm,height=5cm,
  tick label style={font=\footnotesize},label style={font=\small},grid=major,grid style={gray!25},table/col sep=comma,
  x dir=reverse,xlabel={seconds to the cross}]
\nextgroupplot[title={\small indicative price, one session},xmin=0,xmax=300,ymin=99.97,ymax=100.22,
  yticklabel style={/pgf/number format/.cd,fixed,precision=2}]
\addplot[omDef,very thick,const plot] table[x=s,y=price]{figdata/microstructure/10-auction-mechanics-and-theory/path.csv};
\nextgroupplot[title={\small $|$indicative $-$ final$|$ (ticks)},xmode=log,xmin=0.7,xmax=150,ymin=0,ymax=6]
\addplot[omThm,very thick,mark=*,error bars/.cd,y dir=both,y explicit] table[x=s,y=gap,y error=se]{figdata/microstructure/10-auction-mechanics-and-theory/gap.csv};
\end{groupplot}
\end{tikzpicture}

\omcaption{The closing call. Left: the indicative price of one simulated session; the fund's buy pushes it up, responders pull it back, and late orders
set the close. Right: the distance between the last indicative and the final price, against the time left, mean over 24 sessions. Data:
\texttt{mx\_auction.auction\_study}.}\label{fig:mx:auction-mechanics-and-theory:path}
\end{omfigure}

Over 24 simulated sessions (\cref{fig:mx:auction-mechanics-and-theory:path}) the last indicative is 4.7 ticks from the final price 120 seconds before the
cross, 3.5 at 60 seconds, 2.2 at 30 and 1.0 at 10; in the last five seconds, when no more orders arrive, it is the final price. Paired volume arrives
earlier than price information: 76\% of the final paired shares are already paired 120 seconds before the cross, 89\% at 60 and 95\% at 30. The volume is
set early and the price late, because the price is set by the marginal orders.

\begin{dated}{2026-09}{Closing-auction rules on two exchanges}\label{dat:mx:auction-mechanics-and-theory:rules}
Nasdaq publishes closing cross net order imbalance information between 3:50 and 4:00 p.m. ET; its imbalance-only close orders, which exist to offset
imbalances, must be priced and execute only at or above (buy) or below (sell) the 4:00 p.m. bid or ask. On the London Stock Exchange each auction call and
each extension is followed by a random period of up to 30 seconds (the exchange's guide of March 2020).
\end{dated}

\section{Cut-offs, collars and random ends}

\begin{definition}[Auction cut-off time]\label{def:mx:auction-mechanics-and-theory:cutoff}
An \emph{auction cut-off time}\index{auction cut-off time} is the moment after which some orders can no longer be entered, modified or cancelled for the
auction, typically market-on-close and limit-on-close orders, while orders that offset the published imbalance may still be accepted.
\end{definition}

\begin{definition}[Auction collar]\label{def:mx:auction-mechanics-and-theory:collar}
An \emph{auction collar}\index{auction collar} is a price range around a reference price outside which the auction may not clear, or orders may not be
priced; if the uncross would fall outside it, the venue extends the call or clears at the collar.
\end{definition}

\begin{definition}[Random end]\label{def:mx:auction-mechanics-and-theory:randomend}
A \emph{random end}\index{random end} ends an auction call at a random moment within a known window after its scheduled end, so that nobody can be sure
to be the last to enter an order.
\end{definition}

The three answer the same problem: the last order into a call can move its price and nobody can respond. A cut-off stops the orders that would create
imbalances late; a collar bounds how far any order can move the price; a \omterm{def:mx:auction-mechanics-and-theory:randomend}{random end} takes away the certainty of being last. \texttt{firm.auctionsim}
implements the collar as a price band from the start of the call (orders outside it are rejected, \texttt{Out\_J} with reason B), the cut-off with the
engine's freeze control, and the \omterm{def:mx:auction-mechanics-and-theory:randomend}{random end} with the simulator's uniform end time.

\section{Strategic behaviour near the uncross}

What is the last second worth? The experiment adds one more market-on-close buy of 3\,000 shares to each session, sent 60, 10 or 1 second before the
scheduled end, and measures how far it moves the closing price against the same session without it (\cref{fig:mx:auction-mechanics-and-theory:late}).
With a fixed end, the order moves the close by 3.0 ticks (standard error 0.6) sent a minute early, 3.5 ten seconds early and 4.7 (0.4) in the last second:
sent late, it walks the standing interest with no time for responders to offer. With a \omterm{def:mx:auction-mechanics-and-theory:randomend}{random end} of up to 30 seconds, the late order moves the close by
3.7 ticks (0.6): the responders get the rest of the random window. The gain from submitting last, 1.75 ticks with a fixed end, falls to 0.75 with a random
one. That gain is what a participant who wants to move the close (to mark a portfolio, or to trigger something priced off it) buys by waiting; it is also
the price a legitimate late order pays, which is why providers wait too.

\begin{omfigure}
\begin{tikzpicture}
\begin{axis}[width=9cm,height=5.2cm,xmode=log,x dir=reverse,xtick={1,10,60},xticklabels={1,10,60},xlabel={\small seconds before the scheduled end},
  ylabel={\small move of the close (ticks)},tick label style={font=\footnotesize},ymin=0,ymax=6,xmin=0.7,xmax=90,grid=major,grid style={gray!25},
  legend columns=2,legend style={font=\footnotesize,at={(0.5,-0.3)},anchor=north,draw=none,fill=none,
  /tikz/every even column/.append style={column sep=8pt}},table/col sep=comma]
\addplot[omDef,very thick,mark=*,error bars/.cd,y dir=both,y explicit] table[x=tau,y=fixed,y error=fixed_se]{figdata/microstructure/10-auction-mechanics-and-theory/late.csv};
\addplot[omThm,very thick,mark=square*,error bars/.cd,y dir=both,y explicit] table[x=tau,y=random,y error=random_se]{figdata/microstructure/10-auction-mechanics-and-theory/late.csv};
\legend{fixed end,random end (up to 30 s)}
\end{axis}
\end{tikzpicture}

\omcaption{How far a 3\,000-share market-on-close buy moves the closing price, against when it is sent, mean and standard error over 24 simulated
sessions. The last second is worth most with a fixed end. Data: \texttt{mx\_auction.auction\_study}.}\label{fig:mx:auction-mechanics-and-theory:late}
\end{omfigure}

Real closes add what the simulation leaves out: participants who choose their own timing, imbalance-only orders, and large index flows known in advance.
Bogousslavsky and Muravyev (2023) study who trades at the close and what it means for \omterm{def:mx:fragmentation-and-routing:discovery}{price discovery} and liquidity.

\section{Frequent batch auctions}

\begin{definition}[Frequent batch auction]\label{def:mx:auction-mechanics-and-theory:fba}
A \emph{frequent batch auction}\index{frequent batch auction} replaces continuous matching with uniform-price call auctions held at short fixed intervals:
orders arriving within an interval are treated as simultaneous and cleared together at its end.
\end{definition}

Budish, Cramton and Shim (2015) argued that the high-frequency trading arms race is a symptom of continuous matching: when public information moves the
price, a fast trader can take a quote before its owner cancels it, and whoever is fastest wins; they proposed batch auctions, for example every tenth of a
second, so that time becomes discrete and a small speed advantage stops deciding who trades. The simulator runs them with
\texttt{batch\_interval\_ns} (\texttt{firm.auctionsim.batch\_venue}).

The race is easy to stage. A quoter shows one lot a tick either side of a public reference; the reference jumps two ticks 400 times; on each jump the
quoter cancels and requotes after its latency, and a sniper sends an order to take the stale quote after its own. With the quoter 150 microseconds slower
than the sniper, the sniper takes the stale quote on 99.75\% of the jumps in the continuous market, 11.75\% with batches every millisecond, 1.25\% every
10 milliseconds and 0.25\% every 100. With a gap of a millisecond, 99.75\%, 99.75\%, 11.0\% and 0.75\%. The share is close to the latency gap divided by the
batch interval, capped at one: the sniper wins only when a batch closes between its order and the quoter's cancel. A batch interval longer than the
latency gaps that matter turns the race into an auction.

\section{Tutorial: the last ten seconds}\label{tut:mx:auction-mechanics-and-theory:tut}

\begin{tutorial}
\textbf{Goal.} Run closing calls, watch the indicative converge, price the last second under two end rules, and replay the sniping race.
\textbf{End state:} \cref{fig:mx:auction-mechanics-and-theory:path,fig:mx:auction-mechanics-and-theory:late} and the numbers of sections 2, 4 and 5.
\begin{tutsteps}
\item \textbf{Design.} \texttt{firm\_auctionsim.ClosingAuction(call\_s, indicative\_s, cutoff\_s, collar, random\_end\_s)}, then
\texttt{.phases(\dots)} and \texttt{.controls(\dots)} into \texttt{ExchangeConfig} and \texttt{Simulator.add\_events}.
\item \textbf{Sessions.} \texttt{mx\_auction.close\_session(seed, random\_end\_s, extra)}; \texttt{indicatives(res)} and \texttt{final\_cross(res)}.
\item \textbf{Study.} \texttt{auction\_study()}: the gap and arrival paths, and the move of the close by a late order, over 24 seeds.
\item \textbf{Race.} \texttt{snipe\_race(interval\_ms, quoter\_us=\dots, sniper\_us=\dots)} and \texttt{race\_study()}; draw with \texttt{fig\_auction.py}.
\end{tutsteps}
\textbf{What to change next.} Add a cut-off ten seconds before the end for market-on-close orders only, and let responders keep offering; give the
responders a lag of five seconds and watch the \omterm{def:mx:auction-mechanics-and-theory:randomend}{random end}'s protection shrink.
\end{tutorial}

\section{Build: auction policies}\label{bld:mx:auction-mechanics-and-theory:auctionsim}

\begin{build}
\bfield{Purpose} Auction designs for the simulator: the closes of chapters 16 and 28's benchmark algorithms, the halt reopenings of chapter~25, and
Book~11's auction strategies.
\bfield{Interface} \texttt{ClosingAuction(call\_s, indicative\_s, cutoff\_s, collar, random\_end\_s)} with \texttt{.phases}, \texttt{.controls};
\texttt{batch\_venue(interval\_ms)}; \texttt{indicatives(res, venue, locate)}, \texttt{final\_cross(res, venue, locate)}, \texttt{gap\_path(ind, final,
grid\_s)}.
\bfield{Rules} The engine keeps the auction's orders and uncrosses through Book~1's \texttt{firm.auction}; indicatives are published only when some
volume pairs; the collar is a price band from the start of the call; the \omterm{def:mx:auction-mechanics-and-theory:randomend}{random end} is uniform on its window.
\bfield{Acceptance tests} \texttt{code/firm/auctionsim/tests/}: a scripted call whose indicatives and cross are checked by hand; a \omterm{def:mx:auction-mechanics-and-theory:randomend}{random end} inside its
window; an order outside the collar rejected; a batch venue's interval.
\bfield{Stretch} Price-monitoring extensions when the uncross would leave the collar; imbalance-only orders; a closing call on several venues with one
reference price.
\end{build}

\begin{omsources}
\item A.~Madhavan, ``Trading mechanisms in securities markets'', \emph{Journal of Finance} 47(2), 1992.
\item E.~Budish, P.~Cramton and J.~Shim, ``The high-frequency trading arms race: frequent batch auctions as a market design response'', \emph{Quarterly
Journal of Economics} 130(4), 2015.
\item V.~Bogousslavsky and D.~Muravyev, ``Who trades at the close? Implications for price discovery and liquidity'', \emph{Journal of Financial Markets}
66, 2023.
\item Nasdaq, \emph{The Nasdaq Opening and Closing Crosses}, web page, 2026.
\item London Stock Exchange, \emph{Maintaining orderly markets: circuit breakers explained}, 2020.
\end{omsources}

\section{Exercises}

\begin{exercise}[$\star$]\label{exo:mx:auction-mechanics-and-theory:1}
Buys: 300 shares at market, 200 at 10.02, 400 at 10.00. Sells: 500 at 10.00, 300 at 10.01, 200 at 10.03. Find the uncrossing price, the volume and the
surplus with Book~1's rules.
\end{exercise}

\begin{exercise}[$\star$]\label{exo:mx:auction-mechanics-and-theory:2}
The indicative shows 20\,000 shares paired at 50.10 with 3\,000 more to buy. What does that tell a liquidity provider, and at what price would it offer?
\end{exercise}

\begin{exercise}[$\star$]\label{exo:mx:auction-mechanics-and-theory:3}
A sniper is 150 microseconds faster than a quoter. Approximately how often does it take a stale quote under batch auctions every 1 and every 10
milliseconds? And with a gap of a millisecond and batches every 10?
\end{exercise}

\begin{exercise}[$\star\star$]\label{exo:mx:auction-mechanics-and-theory:4}
Standing interest is 500 shares a tick on the sell side. How many ticks does a market-on-close buy of 3\,000 shares move the close if nobody responds? What
did the simulation measure for a last-second order, and why is it less?
\end{exercise}

\begin{exercise}[$\star\star$]\label{exo:mx:auction-mechanics-and-theory:5}
Why is paired volume known long before the price in the simulated call (76\% of the paired shares two minutes before the cross, with the indicative still
4.7 ticks away)?
\end{exercise}

\begin{exercise}[$\star\star$]\label{exo:mx:auction-mechanics-and-theory:6}
Explain why a \omterm{def:mx:auction-mechanics-and-theory:randomend}{random end} reduces the gain from submitting last, and what it costs other participants.
\end{exercise}

\begin{exercise}[$\star\star\star$]\label{exo:mx:auction-mechanics-and-theory:7}
\emph{Coding.} Rerun \texttt{snipe\_race} with batches every 1, 10 and 100 milliseconds and a latency gap of 150 microseconds, and compare the sniping rate
with the gap divided by the interval. Where does the approximation fail?
\end{exercise}

\begin{exercise}[$\star\star\star$]\label{exo:mx:auction-mechanics-and-theory:8}
\emph{Find the flaw.} ``The closing price is set by the whole market, so no single late order can move it much.''
\end{exercise}

\section{Problem: The Last Ten Seconds}

\begin{problem}[{Weekend problem --- the last ten seconds}]\label{pb:mx:auction-mechanics-and-theory:1}
An exchange asks whether to add a \omterm{def:mx:auction-mechanics-and-theory:randomend}{random end} to its closing auction and whether to replace continuous trading with \omterm{def:mx:auction-mechanics-and-theory:fba}{frequent batch auctions}. Answer with the
simulator.

\medskip\noindent\textbf{Part I --- The call.}
\begin{enumerate}
\item State the uncrossing rules and apply them to exercise~1.
\item What does the engine publish during a call, and when?
\item Describe the simulated call's participants.
\item How do the responders price their offers, and why does that matter?
\end{enumerate}

\medskip\noindent\textbf{Part II --- Convergence.}
\begin{enumerate}[resume]
\item Give the gap between the indicative and the final price at 120, 60, 30 and 10 seconds.
\item Why is it zero in the last five seconds?
\item What share of the paired volume is known two minutes and one minute before the cross?
\item What do Nasdaq and the London Stock Exchange publish or randomise?
\end{enumerate}

\medskip\noindent\textbf{Part III --- The last second.}
\begin{enumerate}[resume]
\item Define a cut-off, a collar and a \omterm{def:mx:auction-mechanics-and-theory:randomend}{random end}.
\item How far does a 3\,000-share buy move the close when sent 60, 10 and 1 seconds before a fixed end?
\item And before a \omterm{def:mx:auction-mechanics-and-theory:randomend}{random end} of up to 30 seconds?
\item What is the gain from submitting last under each rule?
\item Who gains from waiting, and who pays?
\end{enumerate}

\medskip\noindent\textbf{Part IV --- Batches and the verdict.}
\begin{enumerate}[resume]
\item State Budish, Cramton and Shim's argument.
\item Describe the sniping race and give its results for the two latency gaps.
\item Derive the approximation gap divided by interval.
\item What does a batch interval cost traders who are not racing?
\item State the \emph{named result}: the gap between the indicative and the final price against the time to the uncross, and the gain from late submission
under a fixed and a \omterm{def:mx:auction-mechanics-and-theory:randomend}{random end}.
\item Should the exchange add a \omterm{def:mx:auction-mechanics-and-theory:randomend}{random end}? Answer in one sentence.
\item Should it switch to batches, and at what interval?
\end{enumerate}
\end{problem}

\section{Interview questions}

\begin{interviewq}[$\star$ \iqroles{trader}]\label{iq:mx:auction-mechanics-and-theory:1}
How is a closing auction's price determined, and what is an imbalance?
\end{interviewq}

\begin{interviewq}[$\star\star$ \iqroles{trader}]\label{iq:mx:auction-mechanics-and-theory:2}
You must buy 500\,000 shares at the close. When do you send the order, and what do you watch?
\end{interviewq}

\begin{interviewq}[$\star\star$ \iqroles{researcher}]\label{iq:mx:auction-mechanics-and-theory:3}
Why do some exchanges end auctions at a random time? What would you measure to see whether it works?
\end{interviewq}

\begin{interviewq}[$\star\star$ \iqroles{researcher}]\label{iq:mx:auction-mechanics-and-theory:4}
Explain \omterm{def:mx:auction-mechanics-and-theory:fba}{frequent batch auctions} and the problem they address.
\end{interviewq}

\begin{interviewq}[$\star\star$ \iqroles{developer}]\label{iq:mx:auction-mechanics-and-theory:5}
Implement the indicative price of a call auction efficiently as orders arrive. What data structure do you keep?
\end{interviewq}

\begin{interviewq}[$\star\star\star$ \iqroles{trader, researcher}]\label{iq:mx:auction-mechanics-and-theory:6}
A competitor always enters large closing orders in the last second. How would you tell a legitimate strategy from an attempt to move the close?
\end{interviewq}
